\begin{definition}
   Квадратную подматрицу, образованную пересечением базисных строк и базисных столбцов, будем называть базисным минором. 
\end{definition}
В данном случае <<минор>> - некоторый жаргонизм, под которым следует понимать просто подматрицу. Сформулируем и докажем теорему о базисном миноре.
\label{Теорема о базисном миноре.}
\begin{theorem}
    Любую строку матрицы можно представить в виде линейной комбинации строк, содержащих базисный минор.
\end{theorem}
\begin{proof}
    Предположим, что строки, содержащие базисный минор, линейно зависимы, тогда и строки базисного минора ЛЗ, что противоречит его определению. Значит строки, содержащие базисный минор, линейно независимы. Тогда их можно взять в качестве базисных строк. Пользуясь теоремой 2.1 получаем требуемое утверждение.
\end{proof}
\begin{note}
Разумеется аналогичная теорема для столбцов тоже верна.
\end{note}
Следующая теорема, называемая теоремой Кронекера-Капелли, играет существенную роль в теории систем линейных уравнений.
\begin{theorem} [Кронекера-Капелли]
    Приписывание к матрице $A$  столбца $b$ не меняет ранг $A$ тогда и только тогда, когда $b$ - линейная комбинация столбцов матрицы $A$.
\end{theorem}
\begin{proof}
    Разобьем доказательство на две части:
    \begin{itemize}
        \item[$\Leftarrow$] Если $b$ - л.к. столбцов матрицы $A$, то элементарными преобразованиями столбцов мы можем привести расширенную матрицу $ \left( \begin{matrix} \text{\large \textit{A}} & |& b \end{matrix} \right) $ к виду $ \left( \begin{matrix} \text{\large \textit{A}} & |& \vec{0} \end{matrix} \right) $. Так как любой набор столбцов, содержащий нулевой столбец, будет ЛЗ, значит ранг получившейся матрицы совпадает с рангом $A$. 
        \item[$\Rightarrow$] Если $rg \left( \begin{matrix} \text{\large \textit{A}} & |& b \end{matrix} \right) =rg \left( \begin{matrix} \text{\large \textit{A}} \end{matrix} \right)=r$, тогда можно выбрать $r$ ЛНЗ столбцов из матрицы $A$. Тогда они будут являться базисными для $ \left( \begin{matrix} \text{\large \textit{A}} & |& b \end{matrix} \right) $. Значит $b$ будет являться линейной комбинацией столбцов матрицы $A$.
    \end{itemize}
\end{proof}
\begin{note}
    Аналогичное утверждение верно и для операции приписывания строк.
\end{note}
Посмотрим какие свойства, связанные с рангами, сохраняются при операциях с матрицами.
\begin{proposition}
\label{rgleq}
    Если для матриц $A$, $B$ определена матрица $AB$, то $$rg(AB) \leqslant min(rg(A), rg(B))$$
\end{proposition}
\begin{proof}
    Заметим, что $rg(AB)\leqslant rg\left( \begin{matrix} \text{\large \textit{A}} & |& \text{\large \textit{AB}} \end{matrix} \right) $ по утверждению 2.2. Но представив матрицу $A$ в виде $(\vec{a}_1, \dots, \vec{a}_n)$, то становиться ясно, что столбцы $AB$ - линейные комбинации столбцов $A$ с коэффициентами из $B$. Тогда элементарными операциями со столбцами можно привести матрицу $\left( \begin{matrix} \text{\large \textit{A}} & |& \text{\large \textit{AB}} \end{matrix} \right) $ к виду $\left( \begin{matrix} \text{\large \textit{A}} & |& 0\end{matrix} \right) $, а у этой матрицы ранг по теореме 2.5 равен $rg(A)$. Значит  $rg(AB)\leqslant rg(A)$.

    Аналогично показываем, что $rg(AB)\leqslant \left( \begin{matrix} \text{\large \textit{B}} & |& \text{\large \textit{AB}} \end{matrix} \right) = \left( \begin{matrix} \text{\large \textit{B}} & |& 0\end{matrix} \right) = rg(B)$, так как строки $AB$ - линейные комбинации строк $B$ с коэффициентами из $A$. Тогда $rg(AB) \leqslant min(rg(A), rg(B))$.
\end{proof}
\begin{proposition}
    Если матрица $A$ невырожденная и определена матрица $AB$, тогда $$rg(AB)=rg(B)$$
\end{proposition}
\begin{proof}
    $A$ - невырожденная $\Rightarrow \exists \ A^{-1}$. Тогда $rg(B)=rg(A^{-1}AB)$ Но по утверждению 2.3 $rg(A^{-1}(AB)) \leqslant rg(AB)$, а также $rg(AB) \leqslant rg(B)$. Значит $rg(AB)=rg(B)$.
\end{proof}
\begin{note}
    Если вы используете определение не вырожденности из учебника Беклемишева (т.е. невырожденная - строки которой ЛНЗ), то в формулировку теоремы необходимо добавить условие о том, что $A$ - квадратная.
\end{note}
Матрица $A_{m \times n}$ имеет \emph{упрощенный вид}, если нижние $m-r$ строк нулевые, а некоторые $r$
 столбцов являются столбцами единичной матрицы. Если эти $r$ столбцов являются первыми столбцами, то матрица $A$ имеет \emph{простейший вид}.
\[
\begin{pmatrix}
a_{11} & 0 & \cdots & 0 & a_{1{r+1}} & \cdots & a_{1n} \\
0 & a_{22} & 0 & \cdots & a_{1{r+1}} & \cdots & a_{2n} \\
\vdots & \vdots & \ddots & \vdots & \vdots & & \vdots \\
0 & 0 & \cdots & a_{rr} & a_{1{r+1}} & \cdots & a_{rn} \\
0 & 0 & \cdots & 0 & 0 & \cdots & 0 \\
\vdots & \vdots & & \vdots & \vdots & & \vdots \\
0 & 0 & \cdots & 0 & 0 & \cdots & 0
\end{pmatrix}
\longrightarrow rg =r
\]
\begin{theorem}
    Любую матрицу элементарными операциями со строками можно привести к упрощенному виду.
\end{theorem}
\begin{proof}
    Будем использовать метод Гаусса для строк:
    \begin{itemize}
        \item Если матрица нулевая, то она уже приведена.
        \item Если матрица не нулевая, то применяем к ней метод Гаусса. Если в какой-то момент мы получаем нулевую строку, то переставляем ее в крайнее нижнее положение и продолжаем алгоритм. 
    \end{itemize}
\end{proof}
Также отметим, что элементарными операциями только со строками базисного минора можно получить <<упрощенный вид>>, в котором единичная матрица расположена на месте базисного минора.
\begin{theorem}
    Если для матриц $A, B$ определена $A+B$, то $$rg(A+B) \leqslant rg(A)+rg(B)$$
\end{theorem}
\begin{proof}
    Пусть $rg(A)=m, rg(B)=k$. Тогда припишем справа к $A+B$ матрицы $A$ и $B$. Получим $rg(A+B) \leqslant rg\left( \begin{matrix} \text{\large \textit{A+B}} & |& \text{\large \textit{A}} & |&\text{\large \textit{B}} \end{matrix} \right)$. Выделим из $A$ и $B$ по набору ЛНЗ столбцов размерами $m$ и $k$ соответственно. Занулим все остальные столбцы элементарными операциями, тогда получим $rg\left( \begin{matrix} \text{\large \textit{0}} & |& \text{\large \textit{A*}} & |&\text{\large \textit{B*}} \end{matrix} \right) \leqslant m+k$, где $A^*$ и $B^*$ матрицы, состоящие только из выбранных базисных столбцов и полученных нулевых столбцов. Отсюда получаем, что $rg(A+B) \leqslant m+k=rg(A)+rg(B)$.
\end{proof}
\section{Системы линейных уравнений}
\label{Системы линейных уравнений.}
В матричном виде СЛУ имеет вид $A\vec{x}=\vec{b}$. \textit{Напоминание}: если $\vec{b}=0$, то система называется однородной, иначе не однородной; если система имеет хотя бы 1 решение, то она называется совместной, иначе несовместной.

Будем делать элементарные операции со строками расширенной матрицы $\left( \begin{matrix} \text{\large \textit{A}} & |& b \end{matrix} \right)$. Предположим, что система $\left( \begin{matrix} P\end{matrix} \right)$ при элементарных операциях со строками перешла в систему $\left( \begin{matrix} Q\end{matrix} \right)$,  тогда если у системы $\left( \begin{matrix} P\end{matrix} \right)$ было решение $\vec{x}$, то $\vec{x}$ будет решением и для системы $\left( \begin{matrix} Q\end{matrix} \right)$. Но в силу обратимости элементарных операций если у  системы $\left( \begin{matrix} Q\end{matrix} \right)$ было решение $\vec{y}$, то $\vec{y}$ будет решением и для системы $\left( \begin{matrix} P\end{matrix} \right)$. Значит системы $\left( \begin{matrix} P\end{matrix} \right)$ и $\left( \begin{matrix} Q\end{matrix} \right)$ эквивалентны.

Далее разрешим единственную операцию со столбцами - перестановку столбцов (то есть по факту перенумерацию переменных в системе).
\begin{theorem}[Кронекера-Капелли]
\label{Теорема Кронекера-Капелли.}
    Система $A\vec{x}=\vec{b}$ совместна $\Leftrightarrow$ $rg(A)=rg\left( \begin{matrix} \text{\large \textit{A}} & |& b \end{matrix} \right)$.
\end{theorem}
\begin{proof}
    Практически сразу получаем доказательство из теоремы 2.5: 
    \begin{itemize}
        \item[$\Rightarrow$] Если существует решение $\vec{x}$, то представив $A=(\vec{a_1}, \dots, \vec{a_n})$, понимаем, что $b$ - линейная комбинация, а значит его приписывание не меняет ранг матрицы.
        \item[$\Leftarrow$] Если при приписывании $b$ ранг матрицы не изменился, то $b$ - л.к. столбцов матрицы $A$, а значит коэффициенты в этой л.к. образуют решение системы.
    \end{itemize}
\end{proof}
\begin{proposition}
    Система $A\vec{x}=\vec{b}$ несовместна тогда и только тогда, когда в упрощенном виде матрицы $\left( \begin{matrix} \text{\large \textit{A}} & |& b \end{matrix} \right)$ есть строка вида $(0, \dots, 0, 1)$.
\end{proposition}
\begin{proof}
    Разобьем доказательство на две части:
    \begin{itemize}
        \item[$\Rightarrow$] Пусть $n, m$ - размеры $A$. Из предыдущей теоремы получаем, что $rg \left( \begin{matrix} \text{\large \textit{A}} & |& b \end{matrix} \right) > rg(A)$ (в другую сторону знак невозможен, так как $A$ - подматрица расширенной матрицы). А так как ранг увеличился, то после приписывания столбца $b$ размер максимального ЛНЗ набора увеличился. Пусть $a_1, \dots, a_k$ - максимальный ЛНЗ набор строк в расширенной матрице. Тогда в нем существует строка $a_i$, в которой первые $m-1$ элементов образуют подстроку, являющуюся линейной комбинацией подстрок, образованных первыми $m-1$ элементами остальных строк из  $a_1, \dots, a_k$ (так как до приписывания $b$, это набор был ЛЗ). Тогда элементарными операциями приводим строку $a_i$ к виду $(0, \dots, 0, b_i)$, где $b_i \neq 0$. Делением на $b_i$ получаем строку $(0, \dots, 0,1)$.
        \item [$\Leftarrow$] Если в упрощенном виде есть строка $(0, \dots, 0, 1)$, то уравнение $x_1 \cdot 0 + \dots + x_{n-1} \cdot 0+ x_n \cdot 0 = 1$ не имеет решения, а значит и вся система не имеет решения.
    \end{itemize}
\end{proof}
Следующее утверждение носит теоретических характер и пригодится в дальнейшем.
\begin{theorem}[Фредгольма]
\label{Теорема Фредгольма.}
    Пусть дана система $Ax=b$. Рассмотрим сопряженную однородную систему $A^Ty=0$. Тогда $Ax=b$ совместна тогда и только тогда, когда любое решение $y$ системы $A^Ty=0$ удовлетворяет системе $y^Tb=0$.
\end{theorem}
\begin{proof}
    Разобьем доказательство на две части:
    \begin{itemize}
        \item[$\Leftarrow$] Пусть $Ax=b$ - несовместна, тогда в упрощенном виде системы $ \left( \begin{matrix} \text{\large \textit{A}} & |& b \end{matrix} \right) $ есть строка $(0, \dots, 0, 1)$. Значит существует линейная комбинация строк $ \left( \begin{matrix} \text{\large \textit{A}} & |& b \end{matrix} \right) $ равная $(0, \dots, 0, 1)$, коэффициенты в этой л.к. примем за $y^T$ (т.е. $y^T \left( \begin{matrix} \text{\large \textit{A}} & |& b \end{matrix} \right)=(0, \cdots, 0, 1)$). Но тогда $y^TA=0$ и $y^Tb=1$ (если по компонентно рассмотреть линейную комбинацию), противоречие. Значит $Ax=b$ - совместна.
        \item[$\Rightarrow$] Если $x$ - решение системы $Ax=b$, то $y^TAx=y^Tb$. Тогда если $y$ - решение системы $y^TA=0$, то $0 \cdot x = 0 = y^Tb$.
    \end{itemize}
\end{proof}